Nowadays, hardware and software systems are everywhere around us. One way to
ensure their correct functioning is to automatically synthesize them from a
formal specification.  This has two advantages over alternatives such as
testing and model checking.  Namely, the design part of the
program-development process can be completely bypassed and the synthesized
program is correct by construction.

In this work we are be interested in synthesizing \emph{reactive
systems}~\cite{hp84}. These are systems which maintain a continuous
interaction with their environment.  Examples of reactive systems include
communication, network, and multimedia protocols as well as operating systems.
For the specification, we consider \emph{linear temporal logic} (LTL, for
short)~\cite{pnueli77}. LTL allows to naturally specify time dependence
amongst events that make up the formal specification of a system. The
popularity of LTL as a formal specification language extends to, amongst
others, artificial intelligence~\cite{gv16,cm19,gnpw20}, hybrid systems and
control~\cite{bvpyb16}, software engineering~\cite{lpb15}, and
bio-informatics~\cite{abbdfhinprs17}.

The
classical doubly-exponential-time synthesis algorithm can be decomposed
into three steps:
\begin{enumerate*}[label={(\roman*)}]
  \item \emph{compile} the input LTL formula into an automaton of exponential
    size~\cite{vw84},
  \item \emph{determinize} the automaton~\cite{safra88,piterman07} incurring a
    second exponential blowup,
  \item and determine the winner of a \emph{two-player zero-sum game} played
    on the latter automaton~\cite{pr89}.
\end{enumerate*}
By and far, most alternative approaches focus on avoiding the determinization
step of the algorithm while keeping the rest of the algorithm intact.  This
has motivated the development of so-called Safra-less approaches,
e.g.~\cite{kpv06,eks16,ekrs17,tushy17} to cite a few.  Worth mentioning are
the on-the-fly game construction implemented in the Strix tool~\cite{msl18}
and the \emph{antichain}-based on-the-fly (bounded) determinization described
in~\cite{fjr09} and implemented in Acacia+~\cite{bbfjr12}. Both techniques
avoid constructing the doubly-exponential deterministic automaton.  It is
worth highlighting that Acacia+ was not ranked in recent editions of the
\emph{Reactive Synthesis Competition} (SYNTCOMP,\footnote{See
\url{http://www.syntcomp.org/}} for short) since it is no longer maintained
despite remaining one of the main temporal synthesis references for new
advancements in the field (see,
e.g.\cite{ffrt17,ztlpv17,apsec20,lms20,bltv20}).

\paragraph*{Contribution}
We provide a complete rewriting of Acacia in C++ articulated around modularity
and leveraging modern techniques for better performances. These techniques
include compile-time specialization of the algorithms, the use of SIMD
registers to store vectors, and several preprocessing steps, one of them
relying on efficient Binary Decision Diagram (BDD) libraries. For now, we have
focused on the natural decision version of the synthesis problem which takes as
input an LTL formula and asks whether a controller satisfying it (no matter
the behavior of the environment) exists. Acacia-Bonsai was specifically
designed to be able to easily try different antichain data structures as well
preprocessing pipelines.
